\documentclass[11pt]{article}
\usepackage{amsmath}
\usepackage{amssymb}
\usepackage{amsthm}
\usepackage{mathtools}
\usepackage{booktabs}
\usepackage{geometry}
\usepackage{hyperref}

\geometry{margin=1in}

\newcommand{\diag}{\operatorname{diag}}
\newcommand{\rank}{\operatorname{rank}}
\newcommand{\rk}{\operatorname{rank}}
\newcommand{\codim}{\operatorname{codim}}
\newcommand{\im}{\operatorname{im}}
\newcommand{\Ker}{\operatorname{Ker}}
\newcommand{\GL}{\operatorname{GL}}
\newcommand{\Orb}{\operatorname{Orb}}
\newcommand{\Stab}{\operatorname{Stab}}
\newcommand{\Aut}{\operatorname{Aut}}
\newcommand{\Hom}{\operatorname{Hom}}
\newcommand{\Rep}{\operatorname{Rep}}
\newcommand{\End}{\operatorname{End}}
\newcommand{\Ext}{\operatorname{Ext}}
\newcommand{\Mat}{\mathrm{Mat}}
\newcommand{\R}{\mathbb{R}}
\newcommand{\Z}{\mathbb{Z}}
\newcommand{\leanchecked}{\par\smallskip\noindent\emph{Lean verification.} This proposition has been checked in Lean.}

\theoremstyle{definition}
\newtheorem{theorem}{Theorem}[section]
\newtheorem{proposition}[theorem]{Proposition}
\newtheorem{lemma}[theorem]{Lemma}
\newtheorem{corollary}[theorem]{Corollary}
\newtheorem{definition}[theorem]{Definition}
\newtheorem{remark}[theorem]{Remark}
\newtheorem{example}[theorem]{Example}

\title{Combinatorics of Saddle Points in Deep Linear Networks}
\author{Guillaume Corlouer}
\date{}

\begin{document}
\maketitle

\begin{abstract}
  
\end{abstract}

\tableofcontents

\section{Introduction: Informal statement of the main theorem}
Flow of the paper:
The main theorem will be the stratification of the critical loci via the stratification of the cyclic variety and its intersection wit the specialization. We introduce notation, we state the main theorem (without any quiver representation theory simply in terms of the rank patterns). 

\section{Related works}

\section{Setup and formal statement}
The critical loci of rank $r$ critical points in the loss landscape of a deep linear network admits a stratification parametrized by the rank of the critical points, rank patterns and a subset $S$ of singular vectors of the teacher matrix $M$.
\begin{definition}[Determinantal, cyclic, and critical loci]
Define
\[
\Sigma_d^r:=\{\theta\in \Omega:\rk(\mu(\theta))=r\} \quad \text{\textit{Determinantal variety}},
\] 
\[
\Gamma_{d}
:=
\left\{
(\theta,W_0)\in \Omega \times \text{Mat}_{d_0 \times d_L}
\ \middle|\
W_{<\ell}W_{0}W_{>\ell}=0
\text{ for all }1\leq \ell\leq L
\right\}, \quad \text{\textit{Cyclic variety}}
\]
Let $\pi:\Omega \times \text{Mat}_{d_0 \times d_L} \longrightarrow \text{Mat}_{d_0 \times d_L}$ be the natural projection. Define:
\[
\Gamma_{d}^r:=\Gamma_{d}\cap \Sigma_d^r, \qquad \Gamma_{d,S}^{r}:=\Gamma_{d}^r\cap\pi^{-1}(\Sigma_{XY}P_Q) = \left\{
\theta\in \Omega
\ \middle|\
W_{<\ell}\Sigma_{XY}P_QW_{>\ell}=0
\text{ for all }1\leq \ell\leq L
\right\}
\]
Let $S$ be a subset of left-singular vectors of $M$, define the $S$-critical loci:
\[
\mathcal C^S_{d} = \{\theta\in\Omega | \ \nabla_{\theta} L(\theta)=0; \mu(\theta)= P_SW^{\star}\}
\]
\end{definition}
\begin{definition}\label{def:cyclic-rank-pattern}
For $\phi = (\theta,W_0)\in\Gamma_{d}$ we define its \emph{cyclic rank pattern} to be the array
\[
r(\phi)=(r_{ij}(\phi))_{0\le i,j\le L}
\quad\text{where}\quad
r_{ij}(\phi):=\mathrm{rank}\bigl(W_{j}...W_i\bigr) \quad \text{for} \quad \text{and} \quad r_{ii}(\phi) = d_i \text{and} \quad r_{i+1i}(\phi) = \mathrm{rank}\bigl(W_i\bigr)
\]
\end{definition}
\begin{theorem}
Define the strata $\mathcal O^S_{\rho} := \mathcal O_{\rho} \cap \Gamma^r_{d,S}$ where  $\mathcal{O}_{\rho}$ is a $G_d$-orbit of the cyclic variety $\Gamma^r_{d}$.
Assuming that the data is non degenerate and the input covariance matrix is invertible then the loci of critical points of rank $r<r_\text{max}$ admits the following stratification parameterized by subset $S$ of singular vectors of the teacher with $r$ elements and admissible rank patterns $\rho\in R_d^{\mathrm{orb}}$:
\begin{align*}
  \mathcal{C}^r_d & = \bigsqcup_{|S|=r}\mathcal{C}^S_d = \bigsqcup_{|S|=r}\bigsqcup_{\substack{\rho\in R_d^{\mathrm{orb}}\\ \operatorname{rank}(\rho)=\operatorname{rank}(\Sigma_{XY}P_Q)}} \mathcal{O}^S_{\rho} \\
\end{align*}
\end{theorem}
\section{Stratification of the cyclic variety with multiplicity patterns}
Goal: Stratify the cyclic variety with cyclic rank patterns. 


\section{From multiplicity patterns to rank patterns}
\begin{proposition}
Assume whitened input i.e. $\Sigma_X = I_{d_L}$ then the $S$ critical loci satisfies:
\[
\mathcal C_{d}^S=
\left\{
\theta\in \Gamma_{d,S}^r
\ \middle|\
\mu(W)=P_SW^\ast
\right\}
\]
\end{proposition}
\section{Stratification of the critical loci}

\section{Discussion}

\end{document}